\subsection{轨道类型}

设 \(G\) 是在分离拓扑空间 \(E\) 上连续作用的拓扑群。对于每一点 \(x\)（属于 \(E\)），记 \(G_x\) 为 \(x\) 在 \(G\) 中的固定子群，并假设典则映射 \(G / G_x \to Gx\) 是同胚；以下两种情形尤其满足此条件：

\begin{enumerate}
\def\labelenumi{\alph{enumi})}
\item
  G 和 E 的拓扑都是离散的；
\item
  G 在 E 上真作用（General Topology, Chap. III, §4, no. 2, Prop. 4），例如 G 是紧的（General Topology, Chap. III, §4, no. 1, Prop. 2）。
\end{enumerate}

记 \(\mathcal{T}\) 为 G 的闭子群的共轭类集合。对于每个 \(x\in \mathrm{E}\)，称 \(x\) 的轨道类型（有时简称 \(x\) 的类型）为 \(\mathrm{G}_x\) 在 \(\mathcal{T}\) 中的类；同一轨道上的两点具有相同的轨道类型（Algebra, Chap. I, §5, no. 2, Prop. 2）；两个轨道具有相同类型，当且仅当它们作为 G-集合同构（Algebra, Chap. I, §5, no. 5, Th. 1）。对于每个 \(t\in \mathcal{T}\)，记 \(\mathrm{E}_{(t)}\) 为 E 中类型为 \(t\) 的点的集合，即类型为 \(t\) 的轨道之并；这是 E 的稳定子集。对于 \(\mathrm{H}\in t\)，也将 \(\mathrm{E}_{(\mathrm{H})}\) 记作 \(\mathrm{E}_{(t)}\)；例如，\(\mathrm{E}_{(\mathrm{G})}\) 是由 G 固定的点组成的 E 的闭子空间。

在 T 上定义如下预序关系：

\(t \leq t'\Longleftrightarrow\) 存在 \(\mathrm{H} \in t\) 和 \(\mathrm{H}' \in t'\)，使得 \(\mathrm{H} \supset \mathrm{H}'\)。

设 \(\Omega\) 和 \(\Omega'\) 是 G 在 E 上的两个轨道，\(t\) 和 \(t'\) 分别是它们的类型；则 \(t \leq t'\) 当且仅当存在从 \(\Omega'\) 到 \(\Omega\) 的 G-态射（必为满射且连续）。

设 \(x, x'\) 属于 E，\(t, t'\) 是它们的类型；则 \(t \leq t'\) 当且仅当存在 \(a \in G\)，使得 \(aG_{x'}a^{-1} \subset G_x\)。

引理 6. 设 G 是李群。

\begin{enumerate}
\def\labelenumi{\alph{enumi})}
\item
  \(\mathbf{G}\) 的每个递降紧子群序列都是最终稳定的。
\item
  设 H 和 \(H'\) 是 G 的两个紧子群，满足 \(H \subset H'\)，且存在从 \(H'\) 到 H 的（拓扑群）同构。则 \(H = H'\)。
\item
  对于关系 \(t \leq t'\)，集合 \(\mathcal{T}\) 是一个诺特有序集（Theory of Sets, Chap. III, §6, no. 5, text preceding Prop. 7）。
\item
  设 \((\mathrm{H}_{i})_{i\geq1}\) 是 G 的一个递降紧子群序列；这些子群都是 G 的李子群（Chap. III, §8, no. 2, Th. 2）。整数序列 \((\dim\mathrm{H}_{i})_{i\geq1}\) 递降，因而最终稳定，所以存在整数 N，使得子群 \(H_{i}\) 在 \(i\geq N\) 时具有相同的单位元连通分支。于是正整数序列 \((\mathrm{H}_{i}:(\mathrm{H}_{i})_{0})_{i\geq N}\) 递降且最终稳定，因此当 i 充分大时 \(H_{i}=H_{i+1}\)。
\item
  设 \(f\) 是从 \(\mathrm{H}'\) 到 \(\mathrm{H}\) 的同构。序列 \((f^n(\mathrm{H}))_{n\geq 0}\) 是 \(\mathrm{G}\) 的递降紧子群序列，故 \(f^n(\mathrm{H}) = f^{n + 1}(\mathrm{H})\) 对充分大的 \(n\) 成立，由 \(a\) 可知。由于 \(f\) 是同构，这推出 \(f(\mathrm{H}) = \mathrm{H} = f(\mathrm{H}')\)，所以 \(\mathrm{H} = \mathrm{H}'\)。
\item
  设 \(t, t' \in \mathcal{T}\) 满足 \(t \leq t'\) 且 \(t' \leq t\)。则存在 \(\mathrm{H}, \mathrm{H}_1 \in t\) 和 \(\mathrm{H}', \mathrm{H}_1' \in t'\)，使得 \(\mathrm{H} \supset \mathrm{H}'\) 且 \(\mathrm{H}_1 \subset \mathrm{H}_1'\)。设 \(g\) 和 \(g'\) 是 \(G\) 中满足 \(\mathrm{H}_1 = g\mathrm{H}g^{-1}\) 和 \(\mathrm{H}_1' = g'\mathrm{H}'g'^{-1}\) 的两个元素；置 \(u = g'^{-1}g\)。则
\end{enumerate}

\[
u \mathrm{H} u ^ {- 1} \subset \mathrm{H} ^ {\prime} \subset \mathrm{H};
\]

由 \(b\)，这推出 \(u\mathrm{H}u^{-1} = \mathrm{H}\)，所以 \(\mathrm{H}' = \mathrm{H}\) 且 \(t' = t\)。因此，集合 \(\mathcal{T}\) 是有序的，并由 \(a\) 可知它是诺特的。

定理 2. 设 \(\mathbf{G}\) 是在 \(\mathbf{X}\) 上真作用的李群，且作用律 \((g,x)\mapsto gx\) 是 \(\mathbf{C}^r\) 类的。假设 \(\mathbf{X}\) 是仿紧的。

\begin{enumerate}
\def\labelenumi{\alph{enumi})}
\item
  将 X 的每一点赋予其轨道类型的映射具有如下半连续性：设 \(x \in X\)，且 \(t \in T\) 是其轨道类型；存在 x 的一个稳定开邻域 U，使得对于任意 \(u \in U\)，u 的类型满足 \(\geq t\)。
\item
  对于所有 \(t \in \mathcal{T}\)，\(\mathrm{X}_{(t)}\) 是 \(\mathrm{X}\) 的子流形，\(\mathrm{X}_{(t)}\) 上由 \(\mathrm{G}\) 的作用诱导的等价关系是正则的（Differentiable and Analytic Manifolds, Results, 5.9.5），且态射 \(\mathrm{X}_{(t)} \to \mathrm{X}_{(t)} / \mathrm{G}\) 是一个丛。
\item
  假设 X/G 是连通的。则 X 中元素的轨道类型集合有最大元 \(\tau\)；此外，\(\mathrm{X}_{(\tau)}\) 是 X 的稠密开子集，且 \(\mathrm{X}_{(\tau)}/\mathrm{G}\) 是连通的。
\end{enumerate}

设 \(x\) 是 \(\mathrm{X}\) 中的一点，且 \(t \in \mathcal{T}\) 是其类型。为证明 \(a)\) 和 \(b)\)，可以将 \(\mathrm{X}\) 替换为包含 \(x\) 的稳定开集，从而（由命题 6）可假设 \(\mathrm{X}\) 形如 \(\mathrm{G} \times^{\mathrm{H}}\mathrm{W}\)，其中 \(\mathrm{W}\) 是紧子群 \(\mathrm{H}\)（属于 \(\mathrm{G}\)）的有限维解析线性表示空间，点 \(x\) 是典则投影 \(p(e,0)\) 的像，其中 \((e,0) \in \mathrm{G} \times \mathrm{W}\)，该典则投影为 \(p: \mathrm{G} \times \mathrm{W} \to \mathrm{G} \times^{\mathrm{H}}\mathrm{W}\)。若 \(u = p(g,y) \in \mathrm{G} \times^{\mathrm{H}}\mathrm{W}\) 且 \(a \in \mathrm{G}\)，则 \(au = u\) 当且仅当存在 \(h \in \mathrm{H}\)，使得 \((ag,y) = (gh^{-1},hy)\)，也就是 \(a \in g\mathrm{H}_yg^{-1}\)。因此，\(\mathrm{G}_u = g\mathrm{H}_yg^{-1}\)；特别地，\(\mathrm{G}_x = \mathrm{H}\)，所以 \(\mathrm{G}_u\) 与 \(\mathrm{G}_x\) 的一个子群共轭，这证明了 \(u\) 的类型 \(\geq t\)，从而得到 \(a)\)。

此外，u 的类型为 t，当且仅当 \(G_{u}\) 在 G 中与 H 共轭，或等价地，当且仅当 \(H_{y}\) 在 G 中与 H 共轭；由引理 6 b)，这意味着 \(H_{y} = H\)，从而 y 被 H 固定。若 \(W'\) 是 W 中由 H 固定的元素组成的向量子空间，则 \(\mathrm{X}_{(t)}\) 可与 \(G \times {}^{H}W'\) 识别，因而也可与 \(G/H \times W'\) 识别，于是得到 b)。

为证明 c)，注意 X/G 连通这一假设推出 X 是纯有限维的：事实上，对于所有 \(k \geq 0\)，记 \(X_{k}\) 为满足 \(x \in X\) 且 \(\dim_{x} X = k\) 的点所成的集合；则 \(X_{k}\) 在 X 中既开又闭，且在 G 的作用下稳定，所以 X 等于某个 \(X_{k}\)。

现在对 X 的维数作归纳来证明 c)，\(\dim X = 0\) 时断言显然。设 \(\tau\) 是 X 中各点轨道类型中的极大元（由引理 6 c)，这样的元素存在）。我们将证明：

\(c^{\prime})\) 对于 \(\mathrm{X}_{(t)}\) 的每个子集 A，若 A 在 \(\mathrm{X}_{(\tau)}\) 中既开又闭且在 G 的作用下稳定，则 A 在 X 中的闭包 \(\overline{\mathrm{A}}\) 是开的。

该断言推出 \(c\)。首先注意，\(\mathrm{X}_{(\tau)}\) 在 \(\mathrm{X}\) 中是开的，由 \(a\) 可知；断言 \(c'\) 推出 \(\overline{\mathrm{X}}_{(\tau)}\) 在 \(\mathrm{X}\) 中既开又闭，因而等于 \(\mathrm{X}\)，因为它在 \(\mathrm{G}\) 的作用下稳定且 \(\mathrm{X} / \mathrm{G}\) 连通。设 \(\mathrm{A}\) 是 \(\mathrm{X}_{(\tau)}\) 的一个非空开闭子集，在 \(\mathrm{G}\) 的作用下稳定；由 \(c'\)，\(\overline{\mathrm{A}}\) 在 \(\mathrm{X}\) 中既开又闭且在 \(\mathrm{G}\) 的作用下稳定，故等于 \(\mathrm{X}\)；这意味着 \(\mathrm{A}\) 在

\(X_{(\tau)}\) 中稠密，因而等于 \(X_{(\tau)}\)。所以，\(X_{(\tau)}/G\) 的每个非空开闭子集都等于 \(X_{(\tau)}/G\)，这证明了 \(X_{(\tau)}/G\) 是连通的。最后，由于 \(X_{(\tau)}\) 在 X 中稠密，由 a) 可知 X 的每一点类型都满足 \(\leq \tau\)；换言之，\(\tau\) 是 X 中各点轨道类型中的最大元。

现在证明 \(c'\)。可假设 A 非空；设 \(x \in \overline{A}\)。只需证明 \(\overline{A}\) 是 x 的一个邻域。为此，如上可假设 \(X = G \times {}^{H}W\)，其中 H 紧，x 是 \((e,0)\) 的典则像。首先假设 H 在 W 上平凡作用：则 X 可与 \((G/H) \times W\) 识别，且 \(X_{(\tau)}/G = X/G\) 同胚于 W，因而连通；于是 A/G = X/G，所以 A = X。以下假设 H 在 W 上并非平凡作用。在 W 上选取一个紧群 H-不变内积；令 S 为 W 中的单位球面（范数为 1 的向量所成的集合）。注意 S/H 是连通的：事实上，若 \(\dim(W) \geq 2\)，S 连通；若 \(\dim(W) = 1\)，S 是一个两点空间，H 在其上非平凡作用。置 \(Y = G \times {}^{H}S\)；这是 X 的余维 1 闭子流形，在 G 的作用下稳定，且 Y/G 同胚于 S/H，因而连通。因此由归纳假设，Y 有一个极大轨道类型 \(\theta\)，集合 \(Y_{(\theta)}\) 在 Y 中开且稠密，且 \(Y_{(\theta)}/G\) 连通。

考虑 \(R_{+}^{*}\) 在 X 上的作用，该作用由作用律 \((\lambda, (g, w)) \mapsto (g, \lambda w)\) 在 \(R_{+}^{*}\) 对 G × W 的作用下取商诱导。X 中在此作用下共轭的两点具有相同的轨道类型；因此，\(\mathrm{X}_{(\theta)}\) 包含 \(\mathbf{R}_{+}^{*}\mathrm{Y}_{(\theta)}\)，后者是 X 的稠密开子集。但由 a)，\(\mathrm{X}_{(\tau)}\) 是开的，因而与 \(\mathrm{X}_{(\theta)}\) 相交，所以 \(\theta = \tau\)。

另一方面，同胚 \((\lambda,w)\mapsto\lambda w\) 将 \(R_{+}^{*}\times S\) 映到 \(W-\{0\}\)（General Topology, Chap. VI, §2, no. 3, Prop. 3），因而诱导出从 \(R_{+}^{*}\times(S/H)\) 到 \((\mathbf{R}_{+}^{*}\mathrm{S})/\mathrm{H}\) 的同胚，也诱导出从 \(R_{+}^{*}\times(Y/G)\) 到 \((\mathbf{R}_{+}^{*}\mathrm{Y})/\mathrm{G}\) 的同胚，以及从 \(R_{+}^{*}\times(\mathrm{Y}_{(\theta)}/\mathrm{G})\) 到 \((\mathbf{R}_{+}^{*}\mathrm{Y}_{(\theta)})/\mathrm{G}\) 的同胚。因此，\((\mathbf{R}_{+}^{*}\mathrm{Y}_{(\theta)})/\mathrm{G}\) 是连通的，而 \(X_{(\tau)}/\mathrm{G}\) 包含一个连通稠密子集，故自身连通（General Topology, Chap. I, §11, no. 1, Prop. 1）。于是 A 等于 \(X_{(\tau)}\)，从而在 X 中稠密，\(\overline{A}\) 是 x 的一个邻域。定理得证。

采用定理 2 c) 中的记号，称 \(\mathrm{X}_{(\tau)}\) 的点为主点，并称它们的轨道为主轨道。若 x 是 X 中的一点，且 \(G \times G_{x}W\) 是 X 中围绕 x 的轨道的一个线性管，则 x 是主点，当且仅当 \(G_{x}\) 在 W 上平凡作用，也就是说，当且仅当该管形如 \((\mathrm{G}/\mathrm{G}_{x}) \times \mathrm{W}\)。

例。1) 设 G 是连通紧李群，通过内自同构作用于自身。G 中元素 x 的固定子群就是 x 在 G 中的中心化子 \(Z(x)\)；它包含每个包含 x 的极大环面。因此，最大轨道类型 \(\tau\) 是 G 的极大环面的共轭类。开集 \(\mathrm{G}_{(\tau)}\) 是 G 的强正则元集合（\(\S5\)，no. 1, Remark）。假设 G 单连通。则 \(\mathrm{G}_{(\tau)}\) 等于 G 的正则元集合 \(G_{r}\)（\(\S5\)，no. 2, Remark 2）；若 A 是 \(\mathfrak{g} = \mathrm{L}(\mathrm{G})\) 的嘉当子代数 t 的一个单房，则复合映射 \(\pi: A \xrightarrow{\exp} G_{r} \longrightarrow G_{r}/\mathrm{Int}(\mathrm{G})\) 是解析流形的同构。事实上，这是一个同胚（§5, no. 2, Cor. of Prop. 2）；取 \(a \in A\)，置 \(t = \exp a\)，并将 \(T_{t}(\mathrm{G})\) 通过平移 \(\gamma(t)\) 与 g 识别。于是切映射 \(T_{a}(\pi)\) 可与典则嵌入 \(t \to g\) 和商映射 \(g \to g/\mathrm{Im}(\mathrm{Ad} t^{-1} - 1)\) 的复合识别。由于 t 是正则的，\(T_{a}(\pi)\) 是同构，故得到所述结果（Differentiable and Analytic Manifolds, Results, 5.7.8）。

\begin{enumerate}
\def\labelenumi{\arabic{enumi})}
\setcounter{enumi}{1}
\tightlist
\item
  设 E 是实欧几里得仿射空间，H 是 E 的一个超平面集，W 是由关于 H 中超平面的正交反射生成的 E 的位移群。假设 H 在 W 下稳定，且带离散拓扑的群 W 在 E 上真作用。
\end{enumerate}

前述结果可应用于 W 在 E 上的作用。E 中一点 x 的固定子群是由 H 中包含 x 的超平面的反射生成的 W 的子群（Chap. V, §3, no. 3, Prop. 2）。因此，最大轨道类型 \(\tau\) 是子群 \(\{Id_{E}\}\) 的类，而 \(\mathrm{E}_{(\tau)}\) 是 E 的各个房之并。注意，在此情形下覆盖 \(\mathrm{E}_{(\tau)} \to \mathrm{E}_{(\tau)}/\mathrm{W}\) 是平凡的；特别地，若 H 非空，则 \(\mathrm{E}_{(\tau)}\) 不连通。

